\documentclass[10pt,a4paper]{amsart}
\usepackage[margin=19mm]{geometry}
\usepackage{amsmath,amssymb,mathtools}
\usepackage[scr=rsfs,cal=euler]{mathalfa}
\usepackage{xfrac}
\usepackage[unicode,hidelinks,bookmarks=false]{hyperref}
\newcommand{\ie}{i.e.\ }
\newcommand{\cf}{cf.\ }
\newcommand{\resp}{respectively\ }
\newcommand{\Subsection}{\subsection}
\providecommand{\nobreakdash}{\nobreak\hbox{-}}
\setlength{\emergencystretch}{2em}
\multlinegap0pt
\usepackage{amssymb}
\usepackage{float}
\newtheorem{lemma}{Lemma}
\title{Adiabatic perturbation theory for the $F=1$ spinor nonlinear Schr\"odinger equation with nonvanishing boundary conditions}
\author{Vassilios M Rothos}
\date{Source: arXiv:2608.14124v1 (CC BY 4.0)}
\begin{document}
\maketitle

\noindent\emph{Source and licence.} Adiabatic perturbation theory for the $F=1$ spinor nonlinear Schr\"odinger equation with nonvanishing boundary conditions. Version arXiv:2608.14124v1 (CC BY 4.0), distributed by arXiv under the Creative Commons Attribution 4.0 International licence, http://creativecommons.org/licenses/by/4.0/, as recorded in the arXiv licence metadata for this exact version and shown on the abstract page https://arxiv.org/abs/2608.14124v1. Full licence notice: \texttt{licenses/arxiv-2608.14124-CC-BY-4.0.txt}.\par\smallskip

\noindent\textit{Editorial scope.} Counted unit: the article's lemma lem\_5\_1 (source0.tex lines 729--741). The article's complete original proof of it is delivered with it (source0.tex lines 744--806, one \texttt{proof} environment of 63 lines). No mathematics is written, repaired or re-derived: the statement and the proof are the article's own lines, in the article's own notation, and no retained line is substituted or re-flowed.  Retained uncounted as the source-local context the proof needs: lines 715--718.  Nothing was omitted from any retained span, so the package carries no omission record.  No retained line contains a citation, so the wrapper carries no bibliography and no bibliography entry.  No retained line contains a figure environment, an image-inclusion command or drawing code, so no image is delivered and the package declares no asset dependency.  Editorial formatting only, in addition: an AMS \texttt{amsart} preamble that restates the article's own declarations and macros used by the retained lines, namely the environments \texttt{lemma} on separate counters, together with the standard packages those lines need.  The retained environments print in this document as Lemma 1 in order of appearance, whereas the article prints the same items with the numbering of its own full document.  The complete native source of the article as distributed in the arXiv e-print for this exact version is bundled unchanged as \texttt{sources/207635/source0.tex}.
\begin{equation}
U_t - V_x + [U,V] = \epsilon P(x,t,k),
\label{eq:ZeroCurvaturePerturbed4}
\end{equation}

\begin{lemma}
\label{lem_5_1}
There exist matrix-valued functions $\mathcal{R}_\pm(x,t,k)$ such that the Jost solutions $\Psi_\pm(x,t,k)$ satisfy
\begin{equation}
\Psi_{\pm,t} = V\Psi_\pm - i\Psi_\pm \Omega(k) + \epsilon \Psi_\pm \mathcal{R}_\pm(x,t,k),
\label{eq:JostTimePerturbed4}
\end{equation}
where $\Omega(k)$ is the diagonal dispersion matrix determined by the asymptotic diagonalization of the $t$-part of the Lax pair. The correction terms $\mathcal{R}_\pm$ admit the representation
\begin{equation}
\mathcal{R}_\pm(x,t,k)=\int_{\pm\infty}^{x}\Psi_\pm^{-1}(y,t,k)\,P(y,t,k)\,\Psi_\pm(y,t,k)\,dy.
\label{eq:Rpm4}
\end{equation}
\end{lemma}

\begin{proof}
The result follows by comparing the perturbed and unperturbed time evolutions of the Jost solutions. 
Since $\Psi_\pm$ solve the $x$-part of the Lax pair exactly and retain their prescribed asymptotic behavior, the effect of the perturbation enters solely through the $t$-evolution. 
This induces a correction term which can be represented in integral form via a variation-of-constants argument, leading to \eqref{eq:Rpm4}.

To make this precise, we introduce the ansatz
\begin{equation}
\Psi_{\pm,t}=V\Psi_\pm - i\Psi_\pm \Omega + \epsilon \Psi_\pm \mathcal{R}_\pm,
\label{eq:JostTimeAnsatz4}
\end{equation}
where $\mathcal{R}_\pm$ is to be determined. The term $-i\Psi_\pm\Omega$ ensures the correct asymptotic oscillatory time dependence of the Jost solutions as $x\to\pm\infty$. 
Since the $x$-equation remains unchanged, consistency of \eqref{eq:JostTimeAnsatz4} with the perturbed zero-curvature relation \eqref{eq:ZeroCurvaturePerturbed4} is obtained by comparing the mixed derivatives $\Psi_{\pm,xt}$ and $\Psi_{\pm,tx}$.

Differentiating \eqref{eq:JostTimeAnsatz4} with respect to $x$ and using $\Psi_{\pm,x}=U\Psi_\pm$, we obtain
\begin{align}
\Psi_{\pm,tx}
&=V_x\Psi_\pm + V\Psi_{\pm,x} - i\Psi_{\pm,x}\Omega + \epsilon \Psi_{\pm,x}\mathcal{R}_\pm + \epsilon \Psi_\pm \mathcal{R}_{\pm,x} \notag\\
&=V_x\Psi_\pm + VU\Psi_\pm - iU\Psi_\pm\Omega + \epsilon U\Psi_\pm \mathcal{R}_\pm + \epsilon \Psi_\pm \mathcal{R}_{\pm,x}.
\label{eq:txderivative4_full}
\end{align}
On the other hand, differentiating $\Psi_{\pm,x}=U\Psi_\pm$ with respect to $t$ yields
\begin{align}
\Psi_{\pm,xt}
&=U_t\Psi_\pm + U\Psi_{\pm,t} \notag\\
&=U_t\Psi_\pm + UV\Psi_\pm - iU\Psi_\pm\Omega + \epsilon U\Psi_\pm \mathcal{R}_\pm.
\label{eq:xtderivative4_full}
\end{align}
Subtracting \eqref{eq:txderivative4_full} from \eqref{eq:xtderivative4_full}, we obtain
\begin{equation}
\Psi_{\pm,xt}-\Psi_{\pm,tx}
=
\bigl(U_t - V_x + [U,V]\bigr)\Psi_\pm - \epsilon \Psi_\pm \mathcal{R}_{\pm,x}.
\label{eq:mixeddifference4}
\end{equation}
By the perturbed zero-curvature relation \eqref{eq:ZeroCurvaturePerturbed4}, this becomes
\[
0=\epsilon P\Psi_\pm - \epsilon \Psi_\pm \mathcal{R}_{\pm,x}.
\]
Hence
\begin{equation}
\Psi_\pm \mathcal{R}_{\pm,x}=P\Psi_\pm,
\end{equation}
and after multiplication from the left by $\Psi_\pm^{-1}$ we arrive at
\begin{equation}
\mathcal{R}_{\pm,x} = \Psi_\pm^{-1}P\Psi_\pm.
\label{eq:Rpmx4}
\end{equation}

Integrating \eqref{eq:Rpmx4} with respect to $x$, we obtain
\[
\mathcal{R}_\pm(x,t,k)=\mathcal{R}_{\pm,0}(t,k)+\int_{x_0}^{x}\Psi_\pm^{-1}(y,t,k)\,P(y,t,k)\,\Psi_\pm(y,t,k)\,dy.
\]
The integration constants are fixed by the asymptotic normalization. Since $P(x,t,k)\to 0$ as $x\to\pm\infty$ and the Jost solutions are normalized to the asymptotic plane-wave basis, consistency with the unperturbed asymptotic time dependence requires that
\[
\mathcal{R}_\pm(x,t,k)\to 0
\qquad\text{as }x\to\pm\infty.
\]
This uniquely determines
\[
\mathcal{R}_\pm(x,t,k)=\int_{\pm\infty}^{x}\Psi_\pm^{-1}(y,t,k)\,P(y,t,k)\,\Psi_\pm(y,t,k)\,dy,
\]
which is precisely \eqref{eq:Rpm4}.
\end{proof}

\end{document}
